\documentclass[12pt,a4paper]{article}

\usepackage[utf8]{inputenc}
\usepackage[T1]{fontenc}
\usepackage[polish]{babel}
\usepackage[a4paper,margin=2.5cm]{geometry}
\usepackage{enumitem}
\usepackage{hyperref}
\usepackage{titlesec}
\usepackage{graphicx}
\usepackage{subcaption}
\usepackage{tabularx}
\usepackage{booktabs}
\usepackage{url}
\usepackage{comment}
\usepackage[]{hypcap}
\usepackage{chngcntr}
\usepackage{float}
\usepackage{ragged2e}

\newcommand{\projecttitle}{System ekspertowy do wyboru destynacji podróżniczej}
\newcommand{\authorone}{Aleksandra Wieczorek}
\newcommand{\authortwo}{Wojciech Wójcik}
\newcommand{\projectdate}{Warszawa, marzec 2026}


\begin{document}
\sloppy

\begin{titlepage}
    \centering

    \includegraphics[width=0.8\textwidth]{politechnika.png}\par
    \vspace{2.5cm}

    {\Large Projekt z przedmiotu Systemy Ekspertowe\par}
    \vspace{1.2cm}

    {\Huge \bfseries \projecttitle \par}
    \vspace{2cm}

    {\Large Autorzy:\par}
    \vspace{0.5cm}
    {\large \authorone\par}
    \vspace{0.3cm}
    {\large \authortwo\par}

    \vfill

    {\large \projectdate\par}
\end{titlepage}
\tableofcontents
\newpage

\section{Wprowadzenie}

Wybór miejsca podróży jest dla wielu osób procesem złożonym i wielokryterialnym. Decyzja ta zależy nie tylko od budżetu, terminu czy dostępnych środków transportu, ale również od indywidualnych preferencji, takich jak oczekiwany styl wypoczynku, poziom aktywności, potrzeba bezpieczeństwa czy zainteresowanie określonymi atrakcjami. W praktyce oznacza to konieczność uwzględnienia wielu czynników jednocześnie, które często wzajemnie na siebie wpływają.

Współczesne serwisy turystyczne dostarczają dużej liczby informacji o kierunkach wyjazdów, jednak najczęściej wymagają samodzielnej analizy ofert i porównywania wielu parametrów. W rezultacie proces podejmowania decyzji może być czasochłonny i nieintuicyjny. System ekspertowy może stanowić wsparcie w tym procesie poprzez uporządkowanie wiedzy dziedzinowej oraz automatyczne generowanie rekomendacji na podstawie odpowiedzi udzielanych w trakcie konsultacji.

Celem projektu jest opracowanie systemu ekspertowego wspomagającego wybór destynacji podróżniczej na podstawie preferencji, ograniczeń oraz oczekiwań użytkownika. System ma prowadzić kontekstowy dialog, gromadzić informacje dotyczące planowanego wyjazdu, a następnie na podstawie reguł zapisanych w bazie wiedzy wyznaczać najbardziej odpowiednie kierunki podróży.

\section{Założenia projektowe}

Projektowany system ekspertowy zostanie zrealizowany w języku Prolog. Przyjęto założenie, że system będzie miał budowę modułową, co ułatwi jego rozwój, testowanie oraz analizę działania.

Podstawowym założeniem jest rozdzielenie bazy wiedzy, danych dziedzinowych, mechanizmu sterowania oraz interfejsu użytkownika. Wiedza o destynacjach i regułach decyzyjnych będzie zapisana w bazie wiedzy, natomiast logika dialogu i prezentacji wyników zostanie wydzielona do osobnych modułów.

System będzie składał się z czterech głównych części: interfejsu użytkownika, bazy wiedzy, modułu znajdowania reguł minimalnych oraz modułu wnioskowania przybliżonego i rozmytego. Baza wiedzy będzie obejmowała około 100 reguł produkcji wykorzystywanych w procesie konsultacji.

Istotnym założeniem jest kontekstowy charakter dialogu z użytkownikiem. Oznacza to, że kolejne pytania będą zależały od wcześniej udzielonych odpowiedzi, a system będzie mógł dochodzić do rekomendacji różnymi ścieżkami. Forma odpowiedzi użytkownika ma być dopasowana do treści pytania i może obejmować zarówno wybór jednokrotny, jak i wielokrotny.

Końcowa rekomendacja nie będzie zapisana jako gotowy fakt, lecz będzie wynikiem procesu wnioskowania opartego na regułach oraz informacjach uzyskanych od użytkownika. Dzięki temu system zachowa charakter systemu ekspertowego i będzie mógł uzasadnić otrzymany wynik.

\newpage
\section{Cel systemu}

Zadaniem systemu jest pomoc użytkownikowi w wyborze miejsca podróży odpowiadającego jego potrzebom. System będzie uwzględniał między innymi:
\begin{itemize}
    \item preferowany typ wypoczynku,
    \item budżet podróży,
    \item preferowany klimat,
    \item długość wyjazdu,
    \item środek transportu,
    \item oczekiwany poziom aktywności,
    \item bezpieczeństwo,
    \item liczbę współpodróżujących osób,
    \item preferencje dotyczące natury, miasta, zabytków lub rozrywki.
\end{itemize}

Końcowym wynikiem działania systemu będzie nie tylko wskazanie destynacji, ale też uzasadnie rekomendacji.  
Wynik może mieć postać:
jednej głównej rekomendacji,
kilku alternatyw,
stopnia dopasowania oraz
krótkiego wyjaśnienia, które przesłanki zadecydowały

\section{Charakterystyka problemu}

Wybór destynacji podróżniczej jest problemem złożonym, ponieważ zależy od wielu czynników, które mogą mieć zarówno charakter obiektywny, jak i subiektywny. Różni użytkownicy mogą preferować odmienne style podróżowania, na przykład wypoczynek plażowy, zwiedzanie miast, kontakt z naturą, aktywność sportową lub podróże budżetowe. Dodatkowo niektóre preferencje są nieostre i trudno je opisać wyłącznie w sposób zero-jedynkowy, co uzasadnia wykorzystanie elementów wnioskowania rozmytego.

Problem ten dobrze nadaje się do modelowania w systemie ekspertowym, ponieważ:
\begin{itemize}
    \item możliwe jest sformułowanie licznych reguł produkcji opisujących zależności między cechami podróży a destynacjami,
    \item dialog z użytkownikiem może być kontekstowy i zależny od wcześniej uzyskanych odpowiedzi,
    \item możliwe jest zastosowanie zarówno klasycznego wnioskowania regułowego, jak i wnioskowania przybliżonego oraz rozmytego,
    \item wyniki mogą być prezentowane w postaci rekomendacji o różnym stopniu dopasowania.
\end{itemize}

\section{Zakres wiedzy dziedzinowej}

Baza wiedzy systemu będzie zawierała informacje o destynacjach oraz ich cechach istotnych z punktu widzenia użytkownika. Przykładowe obszary wiedzy obejmują:
\begin{itemize}
    \item klimat i sezonowość,
    \item koszt podróży i pobytu,
    \item dostępność atrakcji turystycznych,
    \item charakter wypoczynku,
    \item bezpieczeństwo,
    \item odpowiedniość dla rodzin, par, grup znajomych lub osób podróżujących samotnie,
    \item dostępność połączeń transportowych,
    \item stopień zatłoczenia,
    \item możliwości aktywnego wypoczynku.
\end{itemize}

W systemie będą reprezentowane zarówno fakty opisujące właściwości destynacji, jak i reguły wyprowadzające rekomendacje na podstawie preferencji użytkownika.

\section{Przykładowe atrybuty i wartości}

W systemie planuje się wykorzystanie następujących grup atrybutów:

\subsection{Atrybuty dotyczące preferencji użytkownika}
\begin{itemize}
    \item \texttt{typ\_wypoczynku}: plażowy, miejski, górski, aktywny, kulturowy, relaksacyjny,
    \item \texttt{budzet}: niski, średni, wysoki,
    \item \texttt{klimat}: ciepły, umiarkowany, chłodny,
    \item \texttt{srodek\_transportu}: samolot, samochód, pociąg, dowolny,
    \item \texttt{dlugosc\_wyjazdu}: weekend, kilka\_dni, tydzien, dlugi\_pobyt,
    \item \texttt{towarzystwo}: solo, para, rodzina, znajomi,
    \item \texttt{poziom\_aktywnosci}: niski, średni, wysoki.
\end{itemize}

\subsection{Atrybuty opisujące destynacje}
\begin{itemize}
    \item \texttt{ma\_plaze},
    \item \texttt{ma\_zabytki},
    \item \texttt{ma\_gory},
    \item \texttt{dobry\_na\_weekend},
    \item \texttt{bezpieczny},
    \item \texttt{tani},
    \item \texttt{odpowiedni\_dla\_rodzin},
    \item \texttt{dobra\_komunikacja},
    \item \texttt{wysoka\_temperatura\_latem}.
\end{itemize}

\section{Przykładowe reguły}

W systemie zostaną zdefiniowane reguły produkcji łączące preferencje użytkownika z właściwościami destynacji. Poniżej przedstawiono siedem przykładowych reguł opartych na wcześniej zdefiniowanych atrybutach.

\begin{itemize}
    \item Jeżeli \texttt{typ\_wypoczynku = plazowy}, \texttt{klimat = cieply} i \texttt{budzet = wysoki}, to należy preferować destynacje o cechach \texttt{ma\_plaze} oraz \texttt{wysoka\_temperatura\_latem}.

    \item Jeżeli \texttt{typ\_wypoczynku = miejski}, \texttt{dlugosc\_wyjazdu = weekend} i \texttt{srodek\_transportu = pociag}, to należy preferować destynacje o cechach \texttt{dobry\_na\_weekend} oraz \texttt{dobra\_komunikacja}.

    \item Jeżeli \texttt{typ\_wypoczynku = kulturowy} i destynacja ma cechę \texttt{ma\_zabytki}, to powinna otrzymać wyższy priorytet w rekomendacji.

    \item Jeżeli \texttt{typ\_wypoczynku = gorski}, \texttt{poziom\_aktywnosci = wysoki} i destynacja ma cechę \texttt{ma\_gory}, to powinna zostać uznana za dobrze dopasowaną.

    \item Jeżeli \texttt{towarzystwo = rodzina}, \texttt{budzet = sredni} i destynacja jest \texttt{odpowiedni\_dla\_rodzin}, to powinna zostać zakwalifikowana do rekomendacji.

    \item Jeżeli \texttt{towarzystwo = solo}, \texttt{typ\_wypoczynku = miejski} oraz destynacja jest \texttt{bezpieczny} i ma cechę \texttt{dobra\_komunikacja}, to może zostać uznana za odpowiednią dla podróży indywidualnej.

    \item Jeżeli \texttt{budzet = niski} i destynacja jest \texttt{tani}, to powinna być preferowana względem droższych alternatyw.
\end{itemize}

Przedstawione reguły mają charakter przykładowy i pokazują sposób powiązania odpowiedzi użytkownika z właściwościami destynacji. W docelowej wersji systemu liczba reguł będzie większa, tak aby możliwe było uzyskiwanie rekomendacji różnymi ścieżkami wnioskowania.

\section{Redukty i ich zastosowanie w projekcie}

W kontekście systemów opartych na regułach oraz metod odkrywania wiedzy istotnym pojęciem są \textbf{redukty} (ang. \textit{reducts}). Redukt można rozumieć jako \textbf{minimalny zbiór atrybutów}, który zachowuje zdolność do podejmowania tej samej decyzji co pełny zbiór cech. Oznacza to, że część informacji może być nadmiarowa, a poprawna rekomendacja może zostać wyznaczona również na podstawie mniejszej liczby odpowiednio dobranych cech.

W praktyce redukty służą do:
\begin{itemize}
    \item usuwania atrybutów zbędnych lub powtarzających podobną informację,
    \item upraszczania reguł decyzyjnych,
    \item ograniczania liczby pytań zadawanych użytkownikowi,
    \item przyspieszenia procesu wnioskowania,
    \item zwiększenia przejrzystości uzasadnienia końcowej rekomendacji.
\end{itemize}

W projektowanym systemie ekspertowym pojęcie reduktu będzie wykorzystywane do wyznaczania \textbf{minimalnych zestawów informacji potrzebnych do wskazania odpowiedniej destynacji podróżniczej}. Oznacza to, że system nie zawsze będzie musiał zadawać wszystkie możliwe pytania. Jeżeli dla uzyskania wiarygodnej rekomendacji wystarczy określony podzbiór odpowiedzi użytkownika, pozostałe pytania mogą zostać pominięte.

Przykładowo, dla części rekomendacji może się okazać, że kluczowe znaczenie mają jedynie takie cechy jak:
\begin{itemize}
    \item typ wypoczynku,
    \item budżet,
    \item długość wyjazdu,
    \item preferowany środek transportu.
\end{itemize}

Jeżeli na podstawie tych atrybutów system potrafi jednoznacznie zawęzić zbiór możliwych destynacji, to pozostałe informacje, na przykład dotyczące poziomu aktywności lub stopnia zatłoczenia, nie będą konieczne w danym przebiegu konsultacji. Taki minimalny zestaw cech będzie właśnie przykładem reduktu.

W naszym projekcie redukty będą wykorzystywane przede wszystkim w dwóch obszarach:
\begin{enumerate}
    \item \textbf{Optymalizacja dialogu z użytkownikiem} \\
    System będzie dążył do zadawania jedynie tych pytań, które są rzeczywiście potrzebne do osiągnięcia rekomendacji. Dzięki temu konsultacja stanie się krótsza i bardziej naturalna.

    \item \textbf{Wyznaczanie reguł minimalnych} \\
    Na podstawie analizy atrybutów system będzie wskazywał minimalne zestawy przesłanek prowadzących do określonego wniosku. Pozwoli to uprościć bazę wiedzy i lepiej uzasadniać, dlaczego dana destynacja została zaproponowana.
\end{enumerate}

Zastosowanie reduktów jest szczególnie istotne w systemie rekomendacyjnym, ponieważ preferencje użytkownika mogą być liczne i częściowo nakładające się. Dzięki analizie reduktów możliwe będzie wykrycie, które informacje są naprawdę decydujące, a które mają jedynie charakter pomocniczy. W efekcie system stanie się bardziej efektywny, bardziej zrozumiały dla użytkownika oraz łatwiejszy do dalszego rozwijania.

W dalszej implementacji moduł odpowiedzialny za znajdowanie reguł minimalnych będzie analizował zależności między odpowiedziami użytkownika a końcowymi rekomendacjami, tak aby wskazać najmniejsze zbiory warunków wystarczające do podjęcia decyzji. Pozwoli to połączyć klasyczne reguły produkcji z bardziej ekonomicznym i inteligentnym sterowaniem procesem konsultacji.

\vspace{2cm}
\section{Algorytm znajdowania minimalnych reguł wyboru}

Minimalna reguła wyboru to taka reguła, która prowadzi do tej samej decyzji co reguła pierwotna, ale nie zawiera żadnych zbędnych warunków. Oznacza to, że usunięcie dowolnego warunku z reguły minimalnej powoduje zmianę decyzji, brak decyzji albo pojawienie się niejednoznaczności.

W projekcie reguła będzie reprezentowana w postaci:

\begin{center}
\texttt{regula(Id, [Warunek1, Warunek2, ..., WarunekN], Decyzja)}
\end{center}

Algorytm dla każdej reguły z bazy wiedzy wykonuje następujące kroki:

\begin{enumerate}
    \item Pobierz regułę \texttt{R} w postaci listy warunków \texttt{W} oraz decyzji \texttt{D}.
    \item Wygeneruj wszystkie niepuste podzbiory listy warunków \texttt{W}.
    \item Posortuj podzbiory rosnąco według liczby warunków.
    \item Dla każdego podzbioru \texttt{P} sprawdź, czy:
    \begin{itemize}
        \item podzbiór \texttt{P} prowadzi do decyzji \texttt{D},
        \item podzbiór \texttt{P} nie prowadzi do żadnej innej decyzji,
        \item żaden właściwy podzbiór \texttt{P} nie spełnia dwóch powyższych warunków.
    \end{itemize}
    \item Jeżeli wszystkie warunki są spełnione, zapisz \texttt{P} jako minimalną regułę wyboru dla decyzji \texttt{D}.
    \item Powtórz procedurę dla wszystkich reguł w bazie wiedzy.
\end{enumerate}

\newpage
Schemat algorytmu:

\begin{verbatim}
minimalne_reguly = []

dla każdej reguły R:
    W = warunki reguły R
    D = decyzja reguły R

    kandydaci = wszystkie niepuste podzbiory W
    kandydaci = sortuj_rosnaco_po_liczbie_warunkow(kandydaci)

    dla każdego P w kandydaci:
        jeżeli decyzje(P) = [D]:
            jeżeli nie istnieje P_min w minimalne_reguly,
               takie że P_min jest podzbiorem P
               oraz decyzja(P_min) = D:
                    dodaj P -> D do minimalne_reguly
\end{verbatim}

Przykład:

\begin{center}
\texttt{regula(r1, [typ=plazowy, klimat=cieply, budzet=wysoki], morze)}
\end{center}

Jeżeli system sprawdzi, że:

\begin{center}
\texttt{[typ=plazowy, klimat=cieply] -> morze}
\end{center}

oraz żaden z podzbiorów jednoelementowych nie wystarcza do jednoznacznego wyboru decyzji \texttt{morze}, to minimalna reguła ma postać:

\begin{center}
\texttt{regula\_minimalna(r1, [typ=plazowy, klimat=cieply], morze)}
\end{center}

\section{Architektura systemu}

Projekt zostanie podzielony na cztery główne komponenty:

\begin{enumerate}[label=\arabic*.]
    \item \textbf{Interfejs użytkownika} \\
    Odpowiada za prowadzenie dialogu z użytkownikiem, zadawanie pytań oraz prezentację wyników końcowych.

    \item \textbf{Baza wiedzy i danych dziedzinowych} \\
    Zawiera fakty opisujące destynacje oraz reguły produkcji wykorzystywane w procesie wnioskowania.

    \item \textbf{Moduł znajdowania reguł minimalnych} \\
    Odpowiada za identyfikację minimalnych zestawów atrybutów potrzebnych do uzyskania określonych wniosków oraz za prezentację minimalnych reguł.

    \item \textbf{Moduł wnioskowania rozmytego i przybliżonego} \\
    Rozszerza klasyczne wnioskowanie o obsługę nieostrych pojęć, takich jak niski budżet, wysoka atrakcyjność czy umiarkowany poziom aktywności.
\end{enumerate}



\section{Podsumowanie}

Przedstawiony projekt dotyczy opracowania systemu ekspertowego wspomagającego wybór destynacji podróżniczej na podstawie preferencji i oczekiwań użytkownika. System będzie wykorzystywał bazę wiedzy, reguły produkcji oraz mechanizmy wnioskowania do generowania rekomendacji najlepiej dopasowanych do zadanych kryteriów.

Przyjęta architektura modułowa oraz rozdzielenie bazy wiedzy, interfejsu użytkownika i mechanizmów wnioskowania mają zapewnić przejrzystość, możliwość rozwoju oraz poprawność działania systemu. Zastosowanie elementów wnioskowania przybliżonego i rozmytego pozwoli dodatkowo uwzględnić nieostry charakter części preferencji użytkownika.

Opracowany system ma stanowić przykład praktycznego zastosowania metod sztucznej inteligencji i systemów ekspertowych w problemie wspomagania decyzji dotyczących planowania podróży.

\end{document}